\section{Teoretisk og faglig ramme}

TræningsMester vurderes ud fra fire faglige spor. Det første spor handler om motivation, frafald og effektivt engagement. Det andet handler om app- og tracker-evidens samt appkvalitet. Det tredje handler om progression og autoregulering i styrketræning. Det fjerde handler om interaktionsdesign, datakvalitet og sikkerhed. Tilsammen afgør sporene, om appen kan forbinde plan, faktisk træning, feedback og dataansvar i samme system.

\subsection{Motivation, frafald og effektivt engagement}

Vedvarende fysisk aktivitet er tæt knyttet til autonom motivation, oplevet kompetence og støtte, der giver brugeren meningsfuld kontrol \cite{teixeira2012_sdt_exercise}. For en træningsapp betyder det, at feedback skal hjælpe brugeren videre frem for kun at markere, om en plan blev fulgt korrekt.

Frafald er et grundvilkår i digitale sundhedsinterventioner, fordi mange brugere stopper, før interventionen kan få virkning \cite{eysenbach2005_law_of_attrition}. Derfor bliver appens håndtering af ændringer i planen relevant. Hvis en afbrudt session eller manglende detaljelog opleves som et nederlag, kan brugerens vej tilbage til træningen blive sværere.

Effektivt digitalt engagement handler om brug, der støtter den ønskede adfærd frem for blot mange klik eller meget data \cite{yardley2016_effective_engagement_dbci}. COM-B og Behaviour Change Technique Taxonomy giver et sprog for capability, opportunity, motivation, selvmonitorering, mål og feedback \cite{michie2011_behaviour_change_wheel,michie2013_bct_taxonomy_v1}. I TræningsMester bruges disse begreber til at vurdere, om brugeren får en næste handling, der stadig føles legitim efter ændringer i træningen.

\subsection{Apps, trackers og appkvalitet}

Systematiske reviews og meta-analyser viser, at smartphoneapps og activity trackers kan øge fysisk aktivitet hos voksne, men effekterne varierer med intervention, målgruppe og brug \cite{laranjo2021_apps_trackers_meta,romeo2019_smartphone_apps_pa_meta}. Wearables kan også støtte fysisk aktivitet, især når feedback og mål er tydelige \cite{brickwood2019_wearable_trackers_meta}. Denne litteratur gør appsporet relevant for TræningsMester, samtidig med at appens effekt skal vurderes i senere bruger- og effektstudier.

Mobile App Rating Scale (MARS) gør appkvalitet til mere end funktionalitet. Engagement, informationskvalitet, æstetik og brugbarhed indgår også \cite{stoyanov2015_mobile_app_rating_scale}. Analyser af fysisk aktivitetsapps viser desuden, at mange apps kun delvist afspejler evidensinformerede adfærdsprincipper \cite{bondaronek2018_quality_pa_apps,kebede2018_evidence_informed_pa_apps,conroy2014_bct_top_ranked_pa_apps}. Derfor vurderes TræningsMester ikke på funktionsmængde, men på om funktionerne hænger sammen med progression, feedback og lavt betjeningsbesvær.

Litteraturen giver to kriterier for TræningsMester. Appen skal have en tydelig adfærdsfunktion, hvor selvmonitorering, mål, feedback og handlingsstøtte hænger sammen. Interaktionen skal samtidig være praktisk nok til, at brugeren kan bruge den under træning.

\subsection{Forskningsstatus og fagligt hul}

Eksisterende app- og trackerforskning giver et solidt grundlag for at se digital feedback og selvmonitorering som relevante for fysisk aktivitet. Meget af litteraturen handler dog om generel aktivitet, skridt, vægt eller brede mHealth-interventioner. Styrketræning kræver en mere detaljeret model for øvelser, sæt, belastning, progression og ændringer i den planlagte træning.

Det faglige hul ligger derfor i koblingen mellem adfærdsstøtte, programlogik, lavt betjeningsbesvær og sikker dataadgang i én styrketræningsapp. Hvis sporene behandles hver for sig, bliver løsningen enten en generel motivationsapp, en teknisk logbog eller en database med brugerflade. TræningsMester samler design og implementering omkring den praktiske test: om motivation, progression, data og adgang stadig hænger sammen, når brugeren træner anderledes end planlagt.

Den faglige værdi ligger i oversættelsen fra litteratur til app. TræningsMester skaber en konkret model for, hvordan kendte principper om motivation, feedback, progression, datakvalitet og adgangskontrol kan bygges ind i et system til realistisk træningsadfærd.

\subsection{Progression og autoregulering i styrketræning}

Styrketræning adskiller sig fra generel aktivitetsregistrering, fordi progression afhænger af øvelser, sæt, gentagelser, belastning, volumen, intensitet og restitution. ACSM beskriver, hvordan belastning, volumen, øvelsesvalg og frekvens reguleres over tid \cite{acsm2009_progression_models}. Nyere resistance-training-anbefalinger fastholder progression og tilpasning som centrale principper \cite{currier2026_acsm_resistance_training}.

Autoregulering gør det relevant at justere træning efter performance, readiness og dagsform \cite{greig2020_autoregulation_resistance_training,larsen2021_autoregulation_systematic_review}. En app bør derfor skelne mellem planlagt workout, faktisk trackerlog, completion-event og næste relevante træningshandling. Tracker/completion-diagrammet belyser denne adskillelse i TræningsMester.

ecofit-studiet viser, at mHealth kan indgå i en resistance-training-nær intervention \cite{plotnikoff2023_ecofit_rct}. Studiet fungerer som kontekstnær inspiration, mens vurderingen af TræningsMester bygger på design, implementering og verifikation.

Progression skaber også et datakrav. Appen skal kunne se forskel på en planlagt øvelse, en udført øvelse og en gennemført workout uden fulde detaljer. Samtidig kan et for detaljeret krav om logging gøre registreringen tung. Den faglige løsning er derfor ikke maksimal dataregistrering, men passende datagranularitet til formålet.

\subsection{Interaktionsdesign, datakvalitet og sikkerhed}

Interaktionsdesign handler om brugerens konkrete møde med systemet. Kursusgrundlaget om brugerundersøgelser, prototyper og evaluering peger på, at designvalg skal vurderes i den kontekst, hvor systemet bruges \cite{hornbaek2024_brugerundersoegelser,hornbaek2024_prototyper,hornbaek2024_evaluering}. I en træningssituation er brugerens opmærksomhed delt mellem krop, udstyr, omgivelser og app. Lav friktion betyder derfor lavt betjeningsbesvær frem for blot et pænt skærmlayout.

Datakvalitet skal vurderes efter formål og kontekst \cite{mohammed2025_five_facets_data_quality}. En trackerlog giver høj detaljeringsgrad, mens et completion-event kan bevare historik og næste handling, når brugeren ikke logger alle sæt. Data er dermed gode til forskellige formål, og appens model skal gøre forskellen tydelig.

Sikkerhed er en designbetingelse. OWASP MASVS beskriver principper for mobil datahåndtering og backendgrænser, mens Supabase RLS-dokumentationen viser, hvordan adgang kan begrænses efter bruger og rolle \cite{owasp2024_masvs,supabase_rls_docs,supabase_securing_api_docs}. GDPR fungerer som juridisk ramme for persondata \cite{europeanparliament2016_gdpr}. I TræningsMester bruges disse kilder til at vurdere adgangsmodellen. Bilagets sikkerhedsdiagram viser den centrale grænse mellem klient, RLS og Edge Functions, jf. Figur~\ref{fig:tm-security-boundary-summary}.

Rammen peger samlet på ét krav. Appen skal gøre den lette handling mulig og samtidig bevare en meningsfuld og sikker datamodel. Det er denne spænding mellem betjeningsbesvær, datakvalitet og adgangskontrol, der gør TræningsMester relevant for sundhedsinformatik.

\subsection{Operationalisering af rammen}

De faglige spor omsættes til fire vurderingsspørgsmål. Får brugeren en reel handlemulighed, når planen ændrer sig? Modelleres progression som mere end en kalender, så plan, faktisk logging, completion og næste handling kan adskilles? Kan brugerfladen bruges midt i træningen uden unødigt registreringsarbejde? Er datakvalitet og adgangskontrol stærke nok til at bære fleksible flows uden at gøre klienten til sikkerhedsautoritet?

\tmdelkonklusion{Teoriens konsekvens}{Et stærkt designvalg skal kunne forklares fagligt, ses i appen og afgrænses metodisk. Hvis et valg kun findes i litteraturen, er det en intention. Hvis det kun findes i koden, er det en funktion. Først når begge dele mødes, bliver valget fagligt relevant.}
